\documentclass[10pt,a4paper]{amsart}
\usepackage[margin=19mm]{geometry}
\usepackage{amsmath,amssymb,mathtools}
\usepackage[scr=rsfs,cal=euler]{mathalfa}
\usepackage{xfrac}
\usepackage[unicode,hidelinks,bookmarks=false]{hyperref}
\newcommand{\ie}{i.e.\ }
\newcommand{\cf}{cf.\ }
\newcommand{\resp}{respectively\ }
\newcommand{\Subsection}{\subsection}
\providecommand{\nobreakdash}{\nobreak\hbox{-}}
\setlength{\emergencystretch}{2em}
\multlinegap0pt
\usepackage{bm}
\usepackage{bbm}
\usepackage{dsfont}
\usepackage{xspace}
\usepackage{cancel}
\usepackage{mathtools}
\usepackage{amsthm}
\makeatletter
\@ifundefined{claimproof}{\newtheorem{claimproof}{Claimproof}}{}
\@ifundefined{theorem}{\newtheorem{theorem}{Theorem}}{}
\@ifundefined{claim}{\newtheorem{claim}[theorem]{Claim}}{}
\@ifundefined{corollary}{\newtheorem{corollary}[theorem]{Corollary}}{}
\makeatother
\makeatletter
\renewcommand{\AA}{\mathbb{A}}
\def\C{\mathcal{C}}
\def\ZZ{\mathbb{Z}}
\expandafter\ifx\csname claimproof\endcsname\relax
\newenvironment{claimproof}[1]{%
\par\vspace{.1in}\noindent\emph{Proof of Claim:}\space#1}{\qedv}
\fi
\expandafter\ifx\csname customthm\endcsname\relax
\newenvironment{customthm}[1]
  {\renewcommand\theinnercustomthm{#1}\innercustomthm}
  {\endinnercustomthm}
\fi
\def\nmdeg{\operatorname{nmdeg}}
\def\qf{\operatorname{qf}}
\def\trdeg{\operatorname{trdeg}}
\makeatother
\title[A transformal transcendence result for algebraic difference ]{A transformal transcendence result for algebraic difference equations}
\author{Moshe Kamensky, Rahim Moosa}
\date{Source: arXiv:2510.16314v1 (CC BY 4.0)}
\begin{document}
\maketitle

\noindent\emph{Source and licence.} A transformal transcendence result for algebraic difference equations. Version arXiv:2510.16314v1 (CC BY 4.0), distributed under the Creative Commons Attribution 4.0 International licence, as recorded in the arXiv licence metadata for this exact version and shown on the abstract page https://arxiv.org/abs/2510.16314v1. Full rights notice: licenses/arxiv-2510.16314-CC-BY-4.0.txt.\par\smallskip

\noindent\textit{Editorial scope.} Counted unit: the Corollary whose source label is \texttt{c3thm} in the article (source0.tex lines 1844--1858, source label \texttt{c3thm}), delivered together with the article's own complete original proof of it (source0.tex lines 1860--1903, one \texttt{proof} environment of 44 lines). No mathematics is written, repaired or re-derived: statement and proof are the article's own text line for line. Required context retained uncounted: thm:nmdeg (lines 1297--1306); d3 (lines 1470--1484). Every definition, display and label of those lines is retained unchanged. Omitted from this package and not redistributed: 1--1296, 1307--1469, 1485--1843, 1859--1859, 1904--2202. Nothing inside a retained excerpt is omitted or altered: every retained body is delivered exactly as the article's lines are written, so no omission receipt of any kind accompanies this package. Citations: the retained excerpts carry no citation command at all, so no bibliography entry of the article is transcribed and no reference is redistributed. Reference adaptation: none was needed and none was made; every reference inside the delivered unit resolves inside it, so no unresolved reference or citation is produced. Printed slips kept literal: no printed word, symbol, hypothesis or number of the counted statement or of its proof was altered. Layout and class: the article's own class is \texttt{amsart} with packages including xstring, amssymb, enumerate, ifdraft, lineno, showlabels; the wrapper is the lane's pinned \texttt{amsart} editorial wrapper at 10pt on A4, which restates the theorem-like environments the retained text needs and adds the article's own macro definitions that the retained text needs (\texttt{\textbackslash AA}, \texttt{\textbackslash C}, \texttt{\textbackslash ZZ}, \texttt{\textbackslash nmdeg}, \texttt{\textbackslash qf}, \texttt{\textbackslash trdeg}), with extra wrapper packages (bm, bbm, dsfont, xspace, cancel, mathtools). The delivered numbering is the unit's own. Assets: no retained span contains a figure environment, a graphics-inclusion command, a PDF-page inclusion, a table or a TikZ picture; no image file is copied into this package, the package declares no asset dependency and no image file is redistributed with it. Licence: version arXiv:2510.16314v1 of this article is distributed under the Creative Commons Attribution 4.0 International licence, as recorded in the arXiv licence metadata for this exact version and shown on the abstract page https://arxiv.org/abs/2510.16314v1. A full rights notice is delivered at licenses/arxiv-2510.16314-CC-BY-4.0.txt. Characters: every retained excerpt is pure ASCII, so the delivered engine, XeLaTeX with its default Latin Modern text font, reports no missing character.
\begin{theorem}
\label{thm:nmdeg}
Suppose $p\in S_{\qf}(k)$ is rational.
\begin{itemize}
\item[(a)]
$\nmdeg(p)\leq\dim(p)+3$.
\item[(b)]
If $k\subseteq\C$ then $\nmdeg(p)\leq 2$.
\end{itemize}
\end{theorem}

\begin{theorem}\label{d3}
Suppose $p\in S_{\qf}(k)$ is rational and nonalgebraic.
\begin{itemize}
\item[(a)]
Suppose $\dim(p)=1$.
If~$p$ is $4$-disintegrated then it is totally disintegrated.
\item[(b)]
 Suppose $\dim(p)>1$ and let
$\ell:=\max\{3, \nmdeg(p)+1\}$.
If~$p$ is $\ell$-disintegrated then it is totally disintegrated.
\item[(c)]
Suppose $k\subset\C$.
If $p$ is $3$-disintegrated then it is totally disintegrated.
\end{itemize}
\end{theorem}

\begin{corollary}\label{c3thm}
Consider the difference-algebraic equation
\begin{equation}
\label{de}
\sigma^n(y)=f\big(y, \sigma(y),\dots,\sigma^{n-1}(y)\big)\tag{E}
\end{equation}
of order $n\geq 1$, where~$y$ is a single variable and~$f\in k(y_0,\dots,y_{n-1})$.
For each $m\geq 1$, consider the following condition on~$(\ref{de})$:
\begin{itemize}
\item[$(C_m)$]
For any pairwise $\sigma$-disjoint solutions $b_1,\dots,b_m$ of~$(\ref{de})$, the transendence degree of $k\langle b_1,\dots,b_m\rangle$ over~$k$ is~$mn$.
\end{itemize}
Then $(C_{n+4})\implies(C_m)$ for all~$m$.
If $k\subseteq\C$ then $(C_3)\implies(C_m)$ for all~$m$.
\end{corollary}

\begin{proof}
We associate to~(\ref{de}) the $\sigma$-rational variety $(\AA^n, \phi)$ where
$$\phi(y_0,\dots,y_{n-1})=\big(y_1,\dots, y_{n-1}, f(y_0,\dots,y_{n-1})\big).$$
Let $p$ be the generic quantifier-free type of $(\AA^n,\phi)$.
Note that $a=(a_0,a_1,\dots,a_{n-1})$ is a realisation of~$p$ if and only if $a_0$ is a solution to~(\ref{de}), $a_\ell=\sigma^\ell(a_0)$ for all $\ell\leq n-1$, and $\trdeg_k(a)=n$.

\begin{claim}
\label{ctod}
If~$(\ref{de})$ satifies~$(C_m)$  then~$p$ is $m$-disintegrated.
\end{claim}
\begin{claimproof}
Let $a_1,\dots,a_m$ be pairwise $\sigma$-disjoint realisations of~$p$, and write each $a_i=(a_{i,0},\dots,a_{i,n-1})$.
Fixing $r\in\ZZ$, note that if $a_{i,0}=\sigma^r(a_{j,0})$ then
$$a_{i,\ell}=\sigma^\ell(a_{i,0})=\sigma^{\ell+r}(a_{j,0})=\sigma^r(a_{j,\ell})$$
for all $\ell=0,\dots,n-1$, so that $a_i=\sigma^r(a_j)$.
Hence $a_{1,0},\dots,a_{m,0}$ are pairwise $\sigma$-disjoint solutions to~(\ref{de}).
Condition $(C_m)$ therefore implies that $k\langle a_{1,0},\dots,a_{m,0}\rangle$ is of transcendence degree $mn$ over~$k$.
But $k\langle a_{1,0},\dots,a_{m,0}\rangle=k(a_1,\dots,a_m)$, proving that $a_1,\dots,a_m$ are independent over~$k$.
(We are using here that for solutions to the rational type~$p$, algebraic independence agrees with independence.)
\end{claimproof}

\begin{claim}
\label{dtoc}
If~$(\ref{de})$ satisfies~$(C_1)$  and~$p$ is $m$-disintegrated then~$(\ref{de})$ satisfies~$(C_m)$.
\end{claim}
\begin{claimproof}
Suppose $b_1,\dots, b_m$ are pairwise $\sigma$-disjoint solutions to~(\ref{de}).
Let $a_i:=(b_i,\sigma(b_i),\dots,\sigma^{n-1}(b_i))$.
Then $a_1,\dots,a_m$ are also pairwise $\sigma$-disjoint.
As $(C_1)$ implies that each $a_i$ is of transcendence degree~$n$ over~$k$, we have that $a_1,\dots, a_m$ are realisations of~$p$.
It follows from $m$-disintegration that $a_1,\dots, a_m$ are independent over~$k$,
so that $k\langle b_1,\dots,b_m\rangle=k(a_1,\dots,a_m)$ is of transcendence degree~$mn$.
\end{claimproof}

Suppose now that~(\ref{de}) satisfies~$(C_{n+4})$.
By Claim~\ref{ctod}, $p$ is $(n+4)$-disintegrated.
We claim that $p$ is totally disintegrated.
Indeed, if $n=1$ then this follows from Theorem~\ref{d3}(a), whereas if $n>1$ then, as Theorem~\ref{thm:nmdeg}(a) tells us that $\nmdeg(p)\leq n+3$, the total disintegration of~$p$ follows from Theorem~\ref{d3}(b).
Claim~\ref{dtoc} then implies that (\ref{de}) satisfies~$(C_m)$ for all~$m$.

Consider, finally, the case when $k\subset\C$.
If~(\ref{de}) satisfies~$(C_3)$ then $p$ is $3$-disintegrated, by Claim~\ref{ctod}, and hence totally disintegrated, by Theorem~\ref{d3}(c).
So, by Claim~\ref{dtoc}, (\ref{de}) satisfies~$(C_m)$ for all~$m$.
\end{proof}

\end{document}
